\section{What current NVIDIA hardware changes}
\label{sec:blackwell}

The 2009 GPU client ran on hardware with no carry-less multiplier at
all: a binary-field product had to be assembled from AND and XOR, one
bit-slice at a time. Three things are different on current NVIDIA
parts, and between them they explain most of the gap between the
rates in Section~\ref{sec:history} and the ones in this paper.

\subsection{A carry-less multiply in PTX}
\label{sec:blackwell-clmad}

PTX ISA~9.3, shipped with CUDA~13.3, added a carry-less multiply-add,
\texttt{clmad.lo.u64} and \texttt{clmad.hi.u64}, which returns the low
and high halves of a $64\times64$-bit product over $\FF_2$. NVIDIA
documents it for \texttt{sm\_80} and later, so it covers Ampere, Ada,
Hopper and Blackwell but not Turing.\art{ecc2k130/NATIVE-CARRYLESS.md}
Our client uses both halves with zero addends and feeds them into a
three-leaf Karatsuba, so a $131$-bit product is six \texttt{CLMAD}
instructions. The native binary for \texttt{sm\_120} contains real
\texttt{CLMAD} opcodes in its SASS, and the build refuses to compile
the native path with anything older than CUDA~13.3 while keeping the
software product for \texttt{sm\_75} slices of a mixed
binary.\art{ecc2k130/include/packed131.h}\art{ecc2k130/benchmarks/ecc2k-clmad/RESULTS.md}

The gain depends on the part. On the RTX~PRO~6000 the native product
lifted the packed walk by $22.4\%$ ($7.110$ to $8.704$~billion
iterations/s) and roughly halved the instructions issued per update,
which is what later let the batch drop to 16 walks per inversion
(Section~\ref{sec:ledger}).\art{ecc2k130/NATIVE-CARRYLESS.md} On Ada
the effect is larger, because the software product it replaces is
more expensive there: $1.81\times$ on the L40S and $1.88\times$ on the
L4.\art{ecc2k130/ADA-L4-L40S.md, §3} On the ECC2K-95 kernel the same
port gave $1.19\times$ on the walk. The T4, at \texttt{sm\_75}, has no
\texttt{clmad} and runs the software product, which is why it sits at
the bottom of every table in this paper.\art{ecc2k130/T4-G4DN.md}

\subsection{The same instruction at very different rates}
\label{sec:blackwell-rates}

Having the instruction is not the same as having it fast. We measured
issue rates with one probe on three parts in one session
(Table~\ref{tab:clmad-dies}).\art{ecc2k130/TWO-CHAINS.md, §6.1}

\begin{table}[t]
\centering
\small
\caption{Issue rates in lanes per SM-clock, one probe, three parts, one
  session (best of three runs per stream). The last column prices one
  \texttt{CLMAD.lo} in \texttt{LOP3} slots, which is what decides
  whether moving work onto the carry-less unit pays.}
\label{tab:clmad-dies}
\begin{tabular}{@{}lrrrrrr@{}}
\toprule
GPU & \texttt{LOP3} & \texttt{CLMAD.lo} & \texttt{CLMAD} lo+hi & \texttt{IMAD.WIDE} & \texttt{POPC} & \texttt{LOP3}/\texttt{CLMAD.lo} \\
\midrule
RTX PRO 6000 (\texttt{sm\_120}) & 63.9 & 1.69 & 1.99 & 31.3 & 16.1 & 37.8 \\
B200 (\texttt{sm\_100}) & 63.7 & 16.9 & 29.1 & 22.9 & 16.0 & 3.78 \\
H100 (\texttt{sm\_90}) & 59.8 & 29.9 & 30.1 & 22.6 & 16.0 & 2.00 \\
\bottomrule
\end{tabular}
\end{table}

The logic pipe, \texttt{POPC} and shared-memory loads run at the same
per-clock rate on all three. The carry-less unit does not. Nsight
attributes \texttt{CLMAD} to the FP64 pipe, and the workstation
Blackwell die executes FP64 at $1/64$ of its FP32 rate, about two lanes
per SM-clock; the data-centre dies run it ten to eighteen times faster
per clock. On the RTX~PRO~6000 one \texttt{CLMAD.lo} costs as much issue
time as $37.8$ logic instructions, on the B200 $3.78$, on the H100
$2.00$.\art{ecc2k130/TWO-CHAINS.md, §6} The RTX~PRO~4500 matches the
6000 to within $1\%$.\art{ecc2k130/ITERATION-FUNCTION.md, §3.1 and §4.5}

This one number reshapes the optimisation problem. On the workstation
die the carry-less unit is scarce, so the right move is to keep
products on it and everything else off it. On the data-centre dies it
is plentiful and the logic pipe binds, so the right move is the
opposite: push reductions, corrections and even basis changes onto
\texttt{CLMAD}. The same source change can therefore win on one
Blackwell part and lose on another. Moving the top-bit correction onto
\texttt{CLMAD} costs $15\%$ on the RTX~PRO~6000 and gains $13.4\%$ on the
B200; a three-limb Karatsuba loses $30\%$ on the 6000 and gains $4.9\%$
on the B200.\art{ecc2k130/TOP-CLMAD.md}\art{ecc2k130/AUTOSWEEP.md, §3.2 and §4}
Kernel tuning does not travel between dies, and each part in this
paper was tuned on its own.

\subsection{Three-input logic and the rest of the integer pipe}
\label{sec:blackwell-logic}

The other instructions that matter are older but just as important.
\texttt{LOP3} evaluates any three-input Boolean function in one
instruction, which is worth $1.65\times$ on the bitsliced product
(Section~\ref{sec:arith-bitsliced}) and makes a full adder two
instructions.\art{ecc2k130/README.md, ``Findings''} On the RTX~PRO~6000,
\texttt{LOP3}, \texttt{IADD3}, the funnel shift \texttt{SHF},
\texttt{PRMT} and \texttt{IMAD} all issue at about 62--63 lanes per
SM-clock, while \texttt{POPC} and \texttt{FLO} run at quarter rate
($15.8$), so each Hamming weight costs about four logic slots; the
static cost model prices them that way.\art{ecc2k130/THROUGHPUT-CEILING.md}\art{ecc2k130/ITERATION-FUNCTION.md, §3.1}
The prime-field walk of Section~\ref{sec:nist} depends instead on the
carry-chain forms \texttt{mad.lo.cc}, \texttt{madc.hi.cc} and
\texttt{addc}, which exist only at 32-bit granularity; writing them by
hand was worth $1.19$--$1.38\times$ there.

Finally, the integer pipe itself got wider. Blackwell's SM lanes are
unified FP32/INT32, where Ada splits them into an FP32-only half and an
FP32/INT32 half, so the note predicted that an Ada SM would retire
integer logic at half the per-clock rate of \texttt{sm\_120}. The L40S
fits that prediction: with the software product it reached $34.4$
million iterations per second per SM, $0.427$ of the RTX~PRO~6000's
$80.4$ on its shipping build, and \texttt{CLMAD} took it to $62.2$
($0.774$).\art{ecc2k130/ADA-L4-L40S.md}
Because the walk is bound by integer logic and the carry-less unit, not
by memory bandwidth or tensor cores, the part that wins is the one with
the most SMs at the highest clock for the money, which on today's
price lists is the RTX~PRO~6000 rather than the B200
(Section~\ref{sec:survey}).
